\SourceLeaf{362}
\noindent If this minimum is to hold for $t'=0$, then
\[
aC=bD
\]
must hold, which gives
\[
n'=\text{Mod.}\,\frac{a}{D}=\text{Mod.}\,\frac{b}{C}
\]
as the value of the minimum. If it is also to become $=1$, we may let $\varepsilon$ be an arbitrary
quantity whose modulus is $=1$ and set
\[
\varepsilon{}C=b,\qquad \varepsilon{}D=a
\]
and
\[
c=\varepsilon{}B,\qquad d=\varepsilon{}A.
\]
If one writes in place of $a,b,c,d$ the quantities $\sqrt{\varepsilon}.a$,
$\sqrt{\varepsilon}.b$, $\sqrt{\varepsilon}.c$, $\sqrt{\varepsilon}.d$,
which does not diminish the generality, then we simply have
\[
c=B,\qquad d=A
\]
or the transformation formula is
\[
t'=\frac{at-b}{Bt+A},\qquad
t=\frac{At'+b}{-Bt'+a}.
\]

5. The simplest transformation formula is
\[
t'=\frac{t-a}{At+1}.
\]

6. The point halfway between the two points $a,b$ is
\[
\frac{\BareRoot{(1+aB)(1+bB)}-\BareRoot{(1+aA)(1+bA)}}
{-A\BareRoot{(1+aB)(1+bB)}+B\BareRoot{(1+aA)(1+bA)}}
\]
or, if the points opposite $a,b$, namely $-\dfrac{1}{A},-\dfrac{1}{B}$,
are denoted by $a',b'$,
\[
\frac{a'\BareRoot{(a-b')(b-b')}-b'\BareRoot{(a-a')(b-a')}}
{\BareRoot{(a-b')(b-b')}-\BareRoot{(a-a')(b-a')}}.
\]
